%==============================================================================
% Chapitre 69 - La montée en pression dans la chambre
%==============================================================================

\section{Introduction}

La combustion de la charge propulsive produit une grande quantité de gaz à haute
température.

Ces gaz sont initialement confinés dans un volume très réduit : la chambre, la
cartouche et les premiers millimètres du canon.

Cette configuration provoque une augmentation très rapide de la pression.

\begin{aretenir}

La pression résulte principalement de la production de gaz dans un volume
initialement très faible.

\end{aretenir}

\section{Origine de la pression}

Les gaz issus de la combustion occupent un volume bien supérieur à celui de la
poudre dont ils proviennent.

Comme ils sont confinés, ils exercent une pression sur toutes les surfaces qui
les entourent.

Cette pression constitue la source d'énergie mécanique qui mettra ensuite le
projectile en mouvement.

\section{Une pression uniforme ?}

Dans une première approximation, la pression est souvent considérée comme
uniforme à l'intérieur de la chambre.

En réalité, de faibles gradients de pression existent localement, notamment
pendant les toutes premières phases de la combustion.

Ces différences deviennent rapidement négligeables devant la pression moyenne.

\section{Action sur les différents éléments}

La pression agit simultanément sur :

\begin{itemize}
    \item le fond de l'étui ;
    \item les parois de la chambre ;
    \item le projectile ;
    \item la culasse.
\end{itemize}

Ces efforts sont équilibrés par la résistance mécanique de l'arme jusqu'au
déplacement du projectile.

\section{Évolution temporelle}

La pression n'augmente pas instantanément.

Elle croît progressivement au fur et à mesure que la combustion produit de
nouveaux gaz.

Lorsque le projectile commence à avancer, le volume disponible augmente, ce qui
modifie l'évolution de la pression.

\section{Influence du volume}

Le volume disponible joue un rôle fondamental.

À quantité de gaz identique, une augmentation du volume tend à diminuer la
pression.

Le déplacement du projectile constitue donc un mécanisme naturel de limitation
de la montée en pression.

\section{Pression maximale}

La pression atteint généralement un maximum peu après le début du mouvement du
projectile.

Elle décroît ensuite progressivement au cours de son déplacement dans le canon.

La valeur maximale dépend notamment :

\begin{itemize}
    \item de la cartouche ;
    \item de la poudre ;
    \item de la masse du projectile ;
    \item de la longueur de la chambre et du canon.
\end{itemize}

\section{Ordres de grandeur}

Les pressions développées dans les armes modernes peuvent atteindre plusieurs
centaines de mégapascals.

Ces valeurs restent compatibles avec les limites de résistance des armes
lorsqu'elles utilisent des munitions conformes à leur conception.

\section{Influence sur le fonctionnement}

La montée en pression conditionne :

\begin{itemize}
    \item la vitesse initiale du projectile ;
    \item les efforts mécaniques subis par l'arme ;
    \item le fonctionnement des armes semi-automatiques et automatiques ;
    \item la précision du tir.
\end{itemize}

\section{Mesure expérimentale}

La pression peut être mesurée à l'aide de capteurs piézoélectriques installés
dans des canons d'essai.

Les résultats permettent notamment de valider les modèles numériques et les
performances des munitions.

\section{Représentation en simulation}

\begin{simulation}

Dans la plupart des simulateurs, la pression est calculée à partir de modèles
de combustion et de l'évolution du volume situé derrière le projectile.

Le calcul est réalisé à chaque pas de temps afin de déterminer les efforts
appliqués au projectile.

\end{simulation}

\section{Vers la mise en mouvement du projectile}

Lorsque la pression devient suffisante pour vaincre les résistances
mécaniques, le projectile commence à se déplacer.

Ce mouvement marque le début de la phase d'accélération dans le canon.

\section{À retenir}

\begin{aretenir}

\begin{itemize}
    \item Les gaz produits par la combustion sont à l'origine de la pression.
    \item La pression agit dans toutes les directions.
    \item Le déplacement du projectile augmente le volume disponible.
    \item L'évolution de la pression résulte d'un équilibre entre production
          des gaz et augmentation du volume.
    \item La pression constitue le moteur de la balistique intérieure.
\end{itemize}

\end{aretenir}
